\section{Related Work}
\label{sec:related_work}

GNNs~\citep{rdl,relbench,relgnn}, graph transformers~\citep{dbformer,relgt,dwivedi2025relational}, and hybrid GNN-LLM models~\citep{relllm} have been proposed for relational deep learning, but these methods are schema-specific and lack transferability across databases. \citet{snowflake2024relational} investigate predictive modeling with LLMs, yet this approach reduces the relational database to a sequence and is constrained by small context windows that are not tailored to relational structure. Tabular foundation models~\citep{tabpfn,tabicl,carte,portal,tabpfnv2} demonstrate the benefits of pretraining and in-context learning, but are confined to single-table settings and cannot capture multi-table relational structure. The most related efforts are KumoRFM~\citep{kumorfm} and Griffin~\citep{griffin}. KumoRFM cannot operate in zero-shot settings, and its technical details remain undisclosed. Griffin aggregates at the row level before GNN-based propagation. In contrast, RT introduces a schema-agnostic, cell-level architecture with Relational Attention masks, enabling pretraining on diverse databases and robust zero-shot transfer.
% For a comprehensive overview of related work, we refer the reader to App.~\ref{appsec:related_work}.
See App.~\ref{appsec:related_work} for an expanded discussion.

% \textbf{Relational deep learning (RDL).} \citet{rdl} introduced an end-to-end framework for predictive modeling on relational databases using neural networks. At its core, RDL represents a database as a relational entity graph: a temporal, heterogeneous graph where each table is a node-type, each row an individual node, and every primary-foreign key relationship an edge. Initial approaches applied heterogeneous graph neural networks directly to these relational entity graphs~\citep{relbench}, and more recently, advanced message-passing approaches have been proposed to enhance the efficiency of GNNs on relational data~\citep{relgnn}. Transformers have emerged as a way to improve upon the GNN message-passing paradigm. \citet{dbformer} and \citet{relgt} propose transformer-based architectures that achieve better performance than GNNs on relational data. An review of RDL architecture can be found in \cite{dwivedi2025relational}. A key limitation, however, is that these architectures are schema-specific, which prevents pretraining and fine-tuning on diverse database structures. Our Relational Transformer, by design, is schema-agnostic, enabling it to learn from and be directly applied to new, unseen database structures. This design principle allows our architecture to demonstrate foundation model-like capabilities, similar to those recently shown in tabular learning.

% \textbf{Tabular foundation models (TFMs).} 
% Recent advancements in tabular foundation models have demonstrated significant promise, exhibiting capabilities such as in-context learning~\citep{tabpfn, tabicl} and efficient fine-tuning~\citep{carte}. 
% These efforts have explored both supervised~\citep{tabpfn} and self-supervised~\citep{portal, carte}  pretraining on real~\citep{carte} or synthetic~\citep{tabpfn, tabicl} data.
% Extending tabular foundation models to relational data is non-trivial, because not only are there multiple tables, but rows in one table are linked to rows in another by foreign-primary key links.
% We take inspiration for the universal cell encoders/decoders from PORTAL~\citep{portal}, which has a similar handling of column names and text/numeric/datetime data types.
% We also take inspiration from the TabPFNv2~\citep{tabpfnv2} transformer architecture, which uses stacked layers of row-wise and then column-wise attention, except that we also have all-pair attention and use attention masks to capture
% foreign-primary key links.


% \textbf{Relational foundation models (RFMs).}
% While tabular foundation models can, in principle, be applied to relational datasets, they fail to account for the rich, multi-table structure of real-world data. To address this limitation, recent works have begun to develop dedicated relational foundation models.
% For instance, \citet{kumorfm} propose a relational foundation model based on graph-transformers and in-context learning, demonstrating both zero-shot and fine-tuning capabilities. However, their solution is not open-sourced, and the exact pretraining procedure has not been released. Separately, \citet{griffin} design a novel architecture pretrained on a mixture of tabular and relational datasets, showing that fine-tuning improves downstream task performance. Their model, however, differs significantly from our own; it first aggregates information within each table and then utilizes graph neural networks to propagate that information between tables. In contrast, our model uses a cell-level representation of the entire database and employs attention masks to directly represent the foreign-key structure. This allows our approach to reason directly over the relational database in its native, cell-based format, offering a more granular and unified understanding.

% \textbf{Pretrained models for graphs.} 
% Pretrained graph learning models have shown strong success in molecular domains. For example, MoleBERT~\citep{xia2023mole} introduces masked atom modeling and triplet-masked contrastive learning to pretrain GNNs for both node-level and graph-level tasks relevant to drug discovery. \citet{beaini2024towards} scale pretraining by curating massive multi-task molecular datasets with billions of labels, showing that combining quantum and biological data improves low-resource tasks.
% % Complementary to these efforts, DiG~\citep{zheng2024predicting} focuses on modeling equilibrium distributions of molecular systems, providing distributional rather than point predictions for protein, ligand, and catalyst-adsorbate systems. 
% Beyond molecules, \cite{huang2023prodigy} introduced a novel pretraining framework that leverages prompt-based graph representations to enable in-context learning on graphs. GraphAny~\citep{zhaofully} develops a zero-shot node classification framework, grounded in linear least-squares principles, that generalizes across graphs with disjoint feature and label spaces by leveraging LinearGNN ensembles and inductive attention. ULTRA~\citep{galkintowards} targets knowledge graph reasoning, learning universal relational representations that transfer zero-shot to unseen knowledge graphs. Finally, GraphGPT~\citep{zhaographgpt} casts graphs as reversible token sequences via Eulerian paths, enabling transformer-based generative pretraining that scales with model size. 

% % Together, these works highlight the promise of foundation-model principles for graph learning. Our approach differs fundamentally by targeting relational databases: we design a schema-agnostic transformer that directly encodes multi-table structures and foreign–primary key dependencies, enabling pretraining and transfer across heterogeneous relational datasets.



% \textbf{Graphs and Large Language Models.}
% A growing line of research investigates how large language models can be adapted to reason over graph-structured and relational data. \cite{fatemitalk,perozzi2024let} explore parameter-efficient encoders that converts graphs into soft prompts for frozen LLMs, showing that performance strongly depends on the choice of graph serialization and structure encoding. \cite{snowflake2024relational} investigate predictive modeling directly on relational databases with LLMs, demonstrating that careful schema-aware prompt design improves over naive text flattening. Building on this idea,~\cite{wu2025large} propose Rel-LLM, a hybrid architecture that combines GNN encoders with LLMs in a retrieval augmented generation framework. These works highlight the potential of combining graph or relational encoders with LLMs, but they are often limited by small context windows and are not tailored to relational databases. In contrast, our proposed Relational Transformer directly encodes multi-table structure via attention masks, offering a fully end-to-end solution that can operate either independently or alongside large language models.
